% =====================================================================
%  IEEE-STYLE CONFERENCE REPORT  (self-contained)
%  MLE-AI-201 Machine Learning -- Final Project
%  Build:   tectonic report.tex    or    pdflatex report.tex (x2)
%  NOTE: run `python -m src.run_all` first so report/figures/*.png exist.
% =====================================================================
\documentclass[10pt,twocolumn,a4paper]{article}

\usepackage[a4paper,top=0.75in,bottom=1.0in,left=0.68in,right=0.68in,columnsep=0.25in]{geometry}
\usepackage{mathptmx}
\usepackage[T1]{fontenc}
\usepackage[utf8]{inputenc}
\usepackage{graphicx}
\usepackage{amsmath,amssymb}
\usepackage{booktabs}
\usepackage{array}
\usepackage{titlesec}
\usepackage{caption}
\usepackage{cite}
\usepackage[safe]{tipa}  % \textschwa for the Azerbaijani letter in column names
\newcommand{\az}[1]{\emph{#1}}
\usepackage[hidelinks]{hyperref}
\usepackage{url}

\renewcommand{\thesection}{\Roman{section}}
\renewcommand{\thesubsection}{\Alph{subsection}}
\titleformat{\section}[block]{\centering\normalsize\scshape}{\thesection.}{0.5em}{}
\titleformat{\subsection}[block]{\normalsize\itshape}{\thesubsection.}{0.5em}{}
\titlespacing*{\section}{0pt}{1.4ex plus 1ex minus .2ex}{1ex}
\titlespacing*{\subsection}{0pt}{1.0ex plus 1ex minus .2ex}{0.6ex}
\newcommand{\abstractkw}[1]{\noindent\textbf{\textit{#1}}}
\captionsetup{font=small,labelsep=period}
\setlength{\parindent}{1em}

% =====================================================================
\begin{document}

% ------------------------- TITLE BLOCK --------------------------------
\twocolumn[
  \begin{@twocolumnfalse}
  \begin{center}
    {\LARGE\bfseries Baku Apartment Prices from the Inside:\\A From-Scratch Decision Tree and Pegasos SVM on bina.az Listings\par}
    \vspace{1.2em}
    % TODO(team): replace with ALL member names, alphabetical by surname.
    {\large
      H\"useynova, T.\quad ---\quad ---\par}
    \vspace{0.5em}
    {\itshape MLE-AI-201 Machine Learning, AI Academy --- Cohort II\par}
    \vspace{1.4em}
  \end{center}
  \begin{quote}
    \small
    \abstractkw{Abstract---}%
    We study 57{,}421 residential sale listings scraped from bina.az and
    frame them as two tasks: regression on $\log(\text{price})$ and binary
    classification of a price tier derived from the training median
    (208{,}000~AZN). Both models are implemented from scratch in NumPy: a
    CART-style decision tree with Gini, entropy and MSE criteria and a
    vectorised threshold-split search, and a soft-margin linear SVM trained
    with the Pegasos stochastic sub-gradient method, extended to non-linear
    decision functions via random Fourier features. Under matched
    hyperparameters our tree reproduces scikit-learn almost exactly (test
    RMSE on log-price $0.2311$ for both; tier F1 $0.9200$ vs.\ $0.9204$),
    and our Pegasos solver matches \texttt{SGDClassifier} with hinge loss
    (F1 $0.8956$ vs.\ $0.8960$; 5-fold CV $0.899\pm0.003$ vs.\
    $0.900\pm0.003$). Error analysis shows both classifiers fail mostly in
    a narrow band around the tier threshold ($\approx$60\% of mistakes
    within $\pm$15\% of it), and the regressor's relative error more than
    doubles for the luxury decile and for peripheral districts. Bonus
    studies add a from-scratch Ridge baseline, a bagged forest of our own
    trees, and k-means market segments that localise the hardest listings.
    Code:
    \href{https://github.com/YOUR-TEAM/bina-ml-final}{\texttt{github.com/YOUR-TEAM/bina-ml-final}}
    (tag \texttt{v1.0-final}). % TODO(team): put the real repo URL here.

    \vspace{0.6em}
    \abstractkw{Index Terms---}decision tree, support-vector machine,
    Pegasos, random Fourier features, regression, classification,
    real-estate price prediction.
  \end{quote}
  \vspace{1.0em}
  \end{@twocolumnfalse}
]

% =====================================================================
\section{Introduction}
Housing listings are a natural testbed for the core supervised-learning
toolkit: prices are heavy-tailed, features are a mix of numeric and
categorical, and the raw dump is genuinely dirty. This project asks
whether implementations written from first principles --- a decision tree
and a support-vector machine in plain NumPy --- can stand next to
scikit-learn on such data when the comparison is fair: same features, same
splits, matched hyperparameters.

Our contributions are: (i) a cleaned, leakage-free modelling dataset built
from the bina.az sale dump, with every cleaning rule documented and
reproducible; (ii) a from-scratch CART-style tree (Gini/entropy/MSE) and a
from-scratch Pegasos SVM~\cite{pegasos} with a random-Fourier-feature
extension~\cite{rff}, both validated against scikit-learn on
fully-determined toy examples before being trusted on real data; (iii) a
matched benchmark with uncertainty estimates and fit/predict timings; and
(iv) an error analysis that explains \emph{where} the remaining error
lives (near the tier threshold, in the luxury tail, and in peripheral
districts). Three bonus studies (from-scratch Ridge, a forest of our own
trees, geospatial k-means) round out the picture.

\section{Background}
\textbf{Decision trees} grow greedily: at each node the split search picks
the feature--threshold pair maximising the impurity decrease
$\mathrm{IG} = I(\text{parent}) - \sum_c \frac{N_c}{N} I(c)$, where $I$ is
Gini $G = 1 - \sum_k p_k^2$ or entropy $H = -\sum_k p_k \log_2 p_k$ for
classification and the node variance for regression. \textbf{Soft-margin
SVMs} minimise a regularised hinge loss; Shalev-Shwartz et
al.~\cite{pegasos} showed the primal can be solved directly by stochastic
sub-gradient descent with step size $\eta_t = 1/(\lambda t)$ --- no
quadratic programming --- with convergence $\tilde O(1/(\lambda T))$.
Rahimi and Recht~\cite{rff} approximate the RBF kernel by an explicit
random cosine feature map, which lets the same linear solver fit
non-linear boundaries in $O(nD)$ time.

\section{Data and Preprocessing}
\textbf{Source.} The Kaggle dump \emph{binaaz-sale-project}
(sehriyarmemmedli): 100{,}775 scraped rows, 51 columns, Azerbaijani
column names, all prices in AZN.

\textbf{Cleaning (every rule is code in \texttt{src/data\_prep.py}).}
(1)~The scraper revisits listings, so we keep the \emph{latest} scrape per
listing URL ($100{,}775 \to 64{,}454$). (2)~We keep the three residential
categories (new build, old stock, house/garden house), because offices,
land plots and garages quote areas in different units and form a different
market ($\to 58{,}452$). (3)~String columns are parsed to numbers:
\az{Sah\textschwa}\ ``145~m\textsuperscript{2}''$\to$\texttt{area\_m2},
\az{M\textschwa rt\textschwa b\textschwa}\ ``7~/~9''$\to$\texttt{floor},
\texttt{building\_floors}, \az{Torpaq sah\textschwa si}\
``1.3~sot''$\to$\texttt{land\_sot}. (4)~Fixed, documented validity rules
--- price in $[10^4, 3{\times}10^6]$~AZN, area in $[15, 800]$~m$^2$,
coordinates inside Azerbaijan's bounding box, $\le$20 rooms, $\le$40
storeys --- remove 342 outlier rows; exact repost duplicates remove 689
more. Final: \textbf{57{,}421 rows}. Deliberately, no quantile-based or
statistical filters are used, so no statistic of any evaluation split
leaks into preprocessing.

\textbf{Leakage.} We drop \texttt{unit\_price} (price per m$^2$) and
\texttt{total\_price} (the target restated) --- both arithmetic functions
of \texttt{price} --- plus identifiers and free text
(\texttt{owner\_name}, \texttt{shop\_name}, \texttt{address},
\texttt{description}, URLs, scrape timestamps) and popularity counters
(\texttt{views}, vip/featured flags) that are consequences of the listed
price rather than property traits.

\textbf{Features (55).} Numerics (area, rooms, floor, building floors,
floor ratio, land area, lat/lng, haversine distance to the Baku centre);
binary flags (repaired, bill of sale, mortgage available, agent-listed,
agency, top/ground floor); one-hot building category; one-hot for the 30
most frequent districts plus \texttt{other}. Missing numerics are encoded
as a constant plus an explicit \texttt{*\_missing} indicator --- an
imputation scheme with no fitted statistics, hence leakage-free by
construction. Categorical encoding is one-hot because both models consume
numeric thresholds; ordinal codes would impose a false order on districts.

\textbf{Tasks, metrics, split.} Task~A: regression on
$\log(\text{price})$ (prices are heavy-tailed; Fig.~\ref{fig:eda}),
reporting RMSE/MAE/$R^2$ on the log scale and, after back-transforming,
in AZN. Task~B: ``premium'' = price above the \emph{training} median,
208{,}000~AZN, reporting precision/recall/F1/ROC-AUC and the confusion
matrix. A seeded 70/15/15 shuffle gives
train/val/test = 40{,}195/8{,}613/8{,}613. Because the tier label is a
median split \emph{defined on the training set}, the classes are balanced
by construction and the balance carries to the other splits (premium
share 0.499/0.505/0.489), which is exactly what stratification would
enforce; the realised balance is reported rather than imposed. All
hyperparameters are chosen on the validation split only; the test set is
touched once per final model.

\begin{table}[t]
  \centering
  \caption{Dataset summary after cleaning.}
  \label{tab:data}
  \begin{tabular}{@{}lr@{}}
    \toprule
    Property & Value \\
    \midrule
    Rows (after cleaning)       & 57{,}421 \\
    Features used               & 55 \\
    Train / Val / Test split    & 40{,}195 / 8{,}613 / 8{,}613 \\
    Tier threshold (train median) & 208{,}000 AZN \\
    Premium share (tr/val/te, \%) & 49.9 / 50.5 / 48.9 \\
    \bottomrule
  \end{tabular}
\end{table}

\begin{figure}[t]
  \centering
  \includegraphics[width=\linewidth]{figures/fig_eda.png}
  \caption{EDA on the training split: (a) raw prices are heavy-tailed;
  (b) $\log(\text{price})$ is near-normal, dashed line = tier threshold;
  (c) area is the strongest single predictor; (d) price grows with rooms.}
  \label{fig:eda}
\end{figure}

\section{Method}
\subsection{Decision tree (from scratch)}
\texttt{src/decision\_tree.py} implements a CART-style binary tree with a
scikit-learn-like \texttt{fit}/\texttt{predict}/\texttt{predict\_proba}
API. The split search is fully vectorised: per feature, samples are
sorted once per node; cumulative class counts (classification) or
cumulative $\sum y, \sum y^2$ (regression) score \emph{all} candidate
thresholds (midpoints between consecutive distinct values) in one pass,
subject to \texttt{min\_samples\_leaf}. The split objective uses
scikit-learn's weighting,
$\frac{N_t}{N}\big(I_t - \frac{N_l}{N_t}I_l - \frac{N_r}{N_t}I_r\big)$,
so \texttt{min\_impurity\_decrease} is directly comparable. Supported
stopping rules: \texttt{max\_depth}, \texttt{min\_samples\_split},
\texttt{min\_samples\_leaf}, \texttt{min\_impurity\_decrease} (a unit
test asserts each one actually changes the tree). Leaves store the mean
(regression) or class counts (classification), from which
\texttt{predict\_proba} returns leaf frequencies; accumulated weighted
impurity decreases give feature importances. A \texttt{max\_features}
option draws a random feature subset per split; the bonus forest uses it.
A sanity test asserts our tree picks the \emph{identical} feature and
threshold as scikit-learn on tiny fully-determined examples.

\subsection{Support-vector machine (from scratch)}
\texttt{src/svm.py} minimises
\begin{equation}
  J(\mathbf{w},b)=\tfrac{\lambda}{2}\lVert\mathbf{w}\rVert^{2}
  +\tfrac{1}{n}\sum_{i=1}^{n}\max\!\big(0,\,1-y_i(\mathbf{w}^{\top}\mathbf{x}_i+b)\big)
\end{equation}
with $y_i \in \{-1,+1\}$ by mini-batch Pegasos: at step $t$, with
$\eta_t = 1/(\lambda t)$ and a random batch $B_t$ ($|B_t|=64$),
$\mathbf{w} \leftarrow (1-\eta_t\lambda)\mathbf{w} + \frac{\eta_t}{|B_t|}
\sum_{i \in B_t, \,\text{margin}<1} y_i \mathbf{x}_i$; the bias $b$ gets
the same averaged sub-gradient step but no shrinkage (it is not
regularised). The objective is tracked during training
(Fig.~\ref{fig:svm}b). Features are standardised with training statistics
only. After training we report $\lVert\mathbf{w}\rVert$, the margin width
$2/\lVert\mathbf{w}\rVert$, and the support vectors
$\{i : y_i(\mathbf{w}^\top\mathbf{x}_i+b)\le 1\}$. For non-linearity we
chose \textbf{random Fourier features} over kernelised Pegasos: with
$n\approx 40$k, keeping per-example dual coefficients makes every
prediction $O(n)$, whereas RFF ($D=300$,
$z(\mathbf{x})=\sqrt{2/D}\cos(W\mathbf{x}+c)$, $W \sim
\mathcal{N}(0, 2\gamma I)$) keeps training and prediction linear-time and
reuses the identical solver unchanged.

\section{Experimental Setup}
Seed 42 everywhere (split, solver, baselines); \texttt{python -m
src.run\_all} reproduces every number and figure from a clean checkout,
running the toy-data sanity tests first. Tuning on validation only:
Task~A sweeps \texttt{max\_depth} $\in \{2..18\}$ then
\texttt{min\_samples\_leaf} $\in \{1..50\}$; Task~B sweeps Gini vs.\
entropy over depths, and for the SVM $\lambda \in
\{10^{-6}..10^{-2}\}$ and RFF $\gamma \in \{0.005..0.1\}$. Chosen
configurations: tree-A (depth 14, leaf 10), tree-B (Gini, depth 14, leaf
5), SVM ($\lambda{=}10^{-3}$, 30k steps, batch 64; RFF $\gamma{=}0.01$).
Baselines use matched settings on identical splits:
\texttt{DecisionTreeRegressor}/\texttt{Classifier} with the same
depth/leaf/criterion, \texttt{SGDClassifier(loss=hinge,
$\alpha{=}\lambda$)} and \texttt{LinearSVC}/\texttt{SVC} with
$C = 1/(n\lambda)$, all with scaling inside a \texttt{Pipeline}. The
$O(n^2)$ \texttt{SVC}-RBF is trained on a documented 15k subsample.
Uncertainty: 5-fold CV over the pooled train+val data, with the tier
threshold and scaler re-fit inside each fold.

\section{Results}

\begin{table}[t]
  \centering
  \caption{Task A --- price regression, test set (log target; AZN metrics
  after back-transform). B8/B9 are bonus models.}
  \label{tab:taskA}
  \small
  \begin{tabular}{@{}lccc@{}}
    \toprule
    Model & RMSE$_{\log}$ & $R^2_{\log}$ & MAE (AZN) \\
    \midrule
    Ours: DecisionTree (mse)   & \textbf{0.2311} & 0.866 & 44{,}089 \\
    sklearn: DTRegressor        & 0.2311 & 0.866 & 44{,}083 \\
    Ours: Ridge, normal eq.\ [B8] & 0.3032 & 0.769 & 63{,}454 \\
    Ours: RandomForest$\times$30 [B9] & 0.2222 & 0.876 & 44{,}612 \\
    sklearn: RandomForest$\times$30  & 0.2151 & 0.884 & --- \\
    \bottomrule
  \end{tabular}
\end{table}

\begin{table}[t]
  \centering
  \caption{Task B --- price-tier classification, test set.}
  \label{tab:taskB}
  \small
  \begin{tabular}{@{}lcccc@{}}
    \toprule
    Model & Acc & Prec & Rec / F1 & AUC \\
    \midrule
    Ours: DecisionTree        & 0.922 & 0.924 & 0.917 / \textbf{0.920} & 0.960 \\
    sklearn: DTClassifier     & 0.922 & 0.924 & 0.917 / 0.920 & 0.959 \\
    Ours: PegasosSVM (lin.)   & 0.898 & 0.901 & 0.890 / \textbf{0.896} & 0.953 \\
    sklearn: SGD (hinge)      & 0.899 & 0.904 & 0.888 / 0.896 & 0.953 \\
    sklearn: LinearSVC        & 0.890 & 0.892 & 0.884 / 0.888 & 0.952 \\
    Ours: Pegasos + RFF       & 0.849 & 0.846 & 0.846 / 0.846 & 0.931 \\
    sklearn: SVC-RBF (15k)    & 0.844 & 0.849 & 0.828 / 0.838 & 0.923 \\
    \bottomrule
  \end{tabular}
\end{table}

\textbf{Our models vs.\ scikit-learn.} Tables~\ref{tab:taskA} and
\ref{tab:taskB} are the headline. With matched hyperparameters the
from-scratch tree and scikit-learn agree to the third--fourth decimal on
both tasks (test RMSE$_{\log}$ 0.2311 both; F1 0.9200 vs.\ 0.9204), and
5-fold CV confirms the agreement is systematic, not a lucky split:
RMSE$_{\log}$ $0.2311\pm0.0019$ for \emph{both} trees, F1
$0.9228\pm0.0025$ (ours) vs.\ $0.9226\pm0.0024$ (sklearn). The residual
gap is tie-breaking: when several splits give equal gain, we take the
first feature in column order while sklearn samples features randomly;
midpoint thresholds also differ in the last float bits. Pegasos lands
within noise of its sklearn twin, \texttt{SGDClassifier(hinge)}: CV F1
$0.8991\pm0.0031$ vs.\ $0.8999\pm0.0031$. The final linear SVM has
$\lVert\mathbf{w}\rVert = 3.83$, margin $2/\lVert\mathbf{w}\rVert =
0.52$, and 11{,}693 of 40{,}195 training points ($29\%$) are support
vectors --- the two tiers genuinely interleave near the threshold.

\textbf{Bias--variance and hyperparameters.} Fig.~\ref{fig:tree}
shows the classic story: training error falls monotonically with depth
while validation error bottoms out at depth 14 for Task A and the
train--val gap widens past it (0.146 vs.\ 0.233 at depth 18).
For the SVM (Fig.~\ref{fig:svm}a), $\lambda$ trades margin for fit:
growing $\lambda$ from $10^{-6}$ to $10^{-2}$ widens the margin
$28\times$ (0.030$\to$0.845) while validation F1 peaks at
$\lambda = 10^{-3}$ --- stronger regularisation beyond that starts to
underfit. The objective decreases as $\tilde O(1/(\lambda t))$ for every
$\lambda$ (Fig.~\ref{fig:svm}b).

\textbf{Linear beats RBF here.} Both non-linear SVMs --- our Pegasos+RFF
and sklearn's SVC-RBF --- lose to the linear models by 4--6 F1 points.
After one-hot encoding, 38 of 55 features are binary and the informative
numeric directions (area, distance to centre) are already near-linearly
related to the log-price tier, so the RBF geometry mostly adds variance.
That the two independent RBF implementations agree (0.846 vs.\ 0.838)
supports this reading.

\textbf{Speed.} Our vectorised tree fits in $\sim$4.7~s vs.\ sklearn's
0.45~s (Cython) --- a $10\times$ gap, not an algorithmic one (both are
$O(F\,N\log N)$ per level). Pegasos fits in 2.2~s; the $O(n^2)$ SVC-RBF
needs 16.6~s on a 15k \emph{subsample} and 11.6~s just to predict, which
is why RFF was the right non-linearity choice at this scale.

\begin{figure}[t]
  \centering
  \includegraphics[width=\linewidth]{figures/fig_tree_depth.png}
  \caption{Task A validation curves. (a) RMSE$_{\log}$ vs.\ depth: the
  train--validation gap widens with depth, validation bottoms at 14.
  (b) \texttt{min\_samples\_leaf} refines the optimum at that depth.}
  \label{fig:tree}
\end{figure}

\begin{figure}[t]
  \centering
  \includegraphics[width=\linewidth]{figures/fig_svm_lambda.png}
  \vspace{2pt}
  \includegraphics[width=0.72\linewidth]{figures/fig_svm_objective.png}
  \caption{Pegasos SVM. Top: validation F1 (a) and margin width (b) vs.\
  $\lambda$ --- regularisation widens the margin monotonically while F1
  peaks at $10^{-3}$. Bottom: the tracked objective $J(\mathbf{w},b)$
  decreases for every $\lambda$.}
  \label{fig:svm}
\end{figure}

\begin{figure}[t]
  \centering
  \includegraphics[width=\linewidth]{figures/fig_importance.png}
  \caption{(a) Tree impurity-decrease importances (Task A): area
  dominates, then geography. (b) Largest-magnitude SVM weights (Task B,
  standardised features): area and centrality push toward premium;
  peripheral districts and missing-floor (house) indicators push toward
  standard.}
  \label{fig:imp}
\end{figure}

\begin{figure}[t]
  \centering
  \includegraphics[width=\linewidth]{figures/fig_confusion_roc.png}
  \caption{Test-set confusion matrices of both from-scratch classifiers
  and their ROC curves.}
  \label{fig:conf}
\end{figure}

\section{Error Analysis}
\textbf{Feature importances and weights} (Fig.~\ref{fig:imp}) agree
across model families: \texttt{area\_m2} carries $\approx$half of the
tree's total impurity decrease, followed by distance-to-centre,
longitude and district indicators; the SVM's largest standardised
weights are the same variables with interpretable signs.

\textbf{Where the classifiers fail.} 60.0\% (tree) and 60.5\% (SVM) of
all test misclassifications are listings whose true price lies within
$\pm$15\% of the 208{,}000~AZN threshold --- the label is nearly
arbitrary there, so the remaining error is largely irreducible given the
median-split framing, not a model deficiency
(Fig.~\ref{fig:conf}).

\textbf{Where the regressor fails} (Fig.~\ref{fig:err}): median relative
error is $\approx$8--10\% through the mid-market but rises to 13\% in the
cheapest decile and 16\% in the most expensive one --- both tails are
priced on attributes we do not observe (dilapidation on one end;
finish, view and prestige on the other). Geographically the error is
worst in garden-house suburbs such as M\textschwa rd\textschwa kan
($\sim$13\% median) where land value and build quality are poorly
captured by our features. The k-means bonus (Fig.~\ref{fig:kmeans}) localises the same
effect geographically: suburban/garden segments reach 15--21\% median
error vs.\ $\sim$8\% in the dense urban core.

\begin{figure}[t]
  \centering
  \includegraphics[width=\linewidth]{figures/fig_error_analysis.png}
  \caption{Task A error analysis on the test set: (a) median relative
  error by price decile --- the luxury tail is hardest; (b) the 12
  largest districts ranked by median relative error.}
  \label{fig:err}
\end{figure}

\section{Bonus Studies}
\textbf{P8 --- Ridge via normal equations.} $\hat{\mathbf w} =
(X^\top X + \lambda I)^{-1} X^\top y$ (bias unpenalised) matches
sklearn's \texttt{Ridge} to $2\times10^{-15}$ RMSE and fits in 28~ms, but
trails the tree badly (RMSE$_{\log}$ 0.303 vs.\ 0.231): price formation
is non-additive (area $\times$ location interactions) and a linear model
cannot express it.

\textbf{P9 --- a forest of our own trees.} Bagging 30 of our trees
(bootstrap + per-split $\sqrt{d}$ feature subsampling, unlimited depth)
drives test RMSE$_{\log}$ from 0.2311 (single tuned tree) to 0.2222,
tracking sklearn's \texttt{RandomForestRegressor} (0.2151) built with the
same settings; Fig.~\ref{fig:rf} shows error falling monotonically as
trees are added. The residual gap to sklearn comes from tie-breaking and
feature-sampling details, consistent with the single-tree comparison.

\textbf{P10 --- market segments.} k-means ($k{=}8$) on coordinates plus a
2-D PCA view of the feature space (Fig.~\ref{fig:kmeans}) yields
interpretable segments (core Baku, coastal belt, garden suburbs) and
sharpens the error analysis above.

\begin{figure}[t]
  \centering
  \includegraphics[width=0.8\linewidth]{figures/fig_rf_curve.png}
  \caption{Bonus P9: test RMSE$_{\log}$ vs.\ ensemble size for a bagged
  forest of our own trees, against the matched sklearn forest and our
  single tree.}
  \label{fig:rf}
\end{figure}

\begin{figure}[t]
  \centering
  \includegraphics[width=\linewidth]{figures/fig_kmeans.png}
  \caption{Bonus P10: (a) k-means segments on lat/lng; (b) regression
  error by segment --- the periphery is hardest; (c) PCA of the feature
  space coloured by segment.}
  \label{fig:kmeans}
\end{figure}

\section{Discussion and Limitations}
The project's main lesson is that the classical algorithms are small:
$\sim$250 lines of NumPy reproduce scikit-learn's accuracy exactly once
the split objective and the sub-gradient schedule are implemented
carefully; the remaining differences are engineering (Cython speed,
random tie-breaking), not statistics. Honest limitations: (i) listed
prices are \emph{asking} prices, not transactions, so all models learn
seller behaviour; (ii) 99.7\% of residential rows are Baku, so nothing
here generalises to other cities; (iii) the median-split tier makes the
classification ceiling explicit --- most remaining errors sit against the
threshold where the label is noise; (iv) free-text descriptions, which
plausibly carry the luxury-tail signal (renovation quality, view,
furnishing), are unused; (v) our RFF non-linearity was tuned over a small
$\gamma$ grid with $D{=}300$ --- a larger map might close part of its gap,
though SVC-RBF's identical failure suggests the boundary is simply
near-linear.

\section{Conclusion}
We built a decision tree and a Pegasos SVM from scratch, matched
scikit-learn under fair conditions on 57k real listings (tree:
RMSE$_{\log}$ 0.2311 = sklearn; SVM: F1 within 0.0004 of
\texttt{SGDClassifier}), and traced the remaining error to the threshold
band, the luxury tail and peripheral geography. Next steps: text features
from descriptions, a quantile-regression forest for calibrated price
intervals, and testing the tier threshold at other percentiles.

\section*{Tools and Acknowledgements}
Dataset: \emph{Bina.az Sale Project} by S.~Memmedli (Kaggle), used under
its stated licence; we neither re-scraped nor redistribute personal data.
Software: NumPy, pandas, Matplotlib, scikit-learn (baselines/bonuses
only), compiled with Tectonic. \textbf{AI-assistance disclosure:} an AI
assistant (Anthropic Claude) was used to draft code scaffolding, the
experiment harness and portions of this report's text; all modelling
decisions, hyperparameter choices and numbers were produced by the
team's own pipeline (\texttt{python -m src.run\_all}), and every line of
the submitted implementation was reviewed, tested and is defended by the
team. % TODO(team): adjust the disclosure to reflect your actual usage.

% =====================================================================
\begin{thebibliography}{9}
\bibitem{pegasos}
S.~Shalev-Shwartz, Y.~Singer, N.~Srebro, and A.~Cotter,
``Pegasos: primal estimated sub-gradient solver for SVM,''
\emph{Mathematical Programming}, vol.~127, no.~1, pp.~3--30, 2011.
\bibitem{rff}
A.~Rahimi and B.~Recht, ``Random features for large-scale kernel
machines,'' in \emph{Advances in Neural Information Processing Systems},
vol.~20, 2007.
\bibitem{cart}
L.~Breiman, J.~Friedman, R.~Olshen, and C.~Stone,
\emph{Classification and Regression Trees}. Wadsworth, 1984.
\end{thebibliography}

\end{document}
